\documentclass[11pt]{article}
\usepackage[T1]{fontenc}
\usepackage{lmodern}
\usepackage[letterpaper,margin=1in]{geometry}
\usepackage{amsmath,amssymb,amsthm,mathtools,microtype}
\usepackage[hidelinks]{hyperref}
\hypersetup{pdftitle={A sourced mixed Einstein--scalar response on a compact Cauchy slab}}
\newtheorem{theorem}{Theorem}
\newtheorem{lemma}[theorem]{Lemma}
\newcommand{\dd}{\,\mathrm d}
\newcommand{\Lie}{\mathcal L}
\newcommand{\Ric}{\operatorname{Ric}}
\newcommand{\tr}{\operatorname{tr}}
\newcommand{\divg}{\operatorname{div}}
\newcommand{\TT}{\mathrm{TT}}
\newcommand{\supp}{\operatorname{supp}}
\numberwithin{equation}{section}
\title{A sourced mixed Einstein--scalar response on a compact Cauchy slab\\
\large The physical cubic form and its marked source contacts}
\author{}
\date{}
\begin{document}
\maketitle
\begin{abstract}
We construct a nonzero connected retarded mixed tensor--scalar
tree coefficient on a genuinely coupled compact Einstein--massless-scalar
background. The preparation and the probe are closed, time-compact
joint metric--scalar currents for the same action. Their nonlinear
completions are defined by geometric constant-mean-curvature slices
and spatial harmonic marks. A differentiated Ward identity shows why
the bare cubic evaluated on retarded and advanced packets is not gauge
invariant: the sourced equations collapse two of its legs. The
second variations of the actual marked preparation and readout
cancel those two terms. A finite-frequency test is nonzero at a
vacuum reference and persists on a branch with nonconstant active
matter. With one positive physical Hadamard covariance, it gives a
normalized coherent one-vertex coefficient. The result treats one
sourced perturbative response; it does not construct a positive
interacting state or a sewing of independently prepared regions.
\end{abstract}

\section{The response to be tested}

Consider the four-dimensional action, in units \(8\pi G=1\),
\begin{equation}
 S[g,X]=\frac12\int_M (R_g-|dX|_g^2)\dd\mu_g
 \label{eq:action}
\end{equation}
with signature \((-+++)\). All action derivatives are relative
derivatives with fixed outer collar data, or equivalently
integrals on the compact temporal overlap when evaluating the
sourced trilinear coefficient below. We use real fields and real
linear sources. At an exact solution \(\bar Y=(\bar g,\bar X)\), put
\(H=S''_{\bar Y}\) and \(C=S'''_{\bar Y}\). We write
\(\langle Hv,w\rangle=H(v,w)\), using the dual pairing of a
field variation and a density. The gauge differential is
\(K\xi=\Lie_\xi\bar Y=(\Lie_\xi\bar g,\xi(\bar X))\).
A current \(J\) is \emph{closed} if \(\langle J,K\xi\rangle=0\)
for every smooth compactly supported \(\xi\).
Use a matched geometric chart to extend \(J\) to an invariant
source functional \(O_J\) with \(O_J'(\bar Y)=J\).
Its Hessian is \(Q_J\).

Our incoming variation is a homogeneous physical Jacobi field \(U\).
A retarded leg \(r\) solves \(Hr=J_1\) and an advanced leg \(a\)
solves \(Ha=J_2\), for two closed currents in separated time
collars. We take both legs in one compatible wave gauge.
The coefficient of interest, with source action \(S-\lambda O_{J_1}\)
and readout \(O_{J_2}\), is
\begin{equation}
 {\cal R}(U;J_1,J_2)
 =-C(U,r,a)+Q_{J_1}(U,a)+Q_{J_2}(U,r).
 \label{eq:response}
\end{equation}
We derive this formula, the actual Hessians, and a nonzero
instance below. Changing the sign of the preparation term
changes the overall sign of the response.

\begin{theorem}\label{thm:main}
There are smooth compact Cauchy slabs of four-dimensional
Einstein--massless-scalar spacetimes with \(d\bar X\ne0\) on the
open witness patch in the tested slab, and a geometric marking such that the following hold.
\begin{enumerate}
 \item One common exact action and one positive physical
 Hadamard representation admit a nonzero incoming physical
 tensor-like Jacobi packet and compact-in-time closed preparation
 and probe currents. The two currents have scalar-like physical
 classes but generally contain metric as well as scalar parts.
 Their sourced retarded coefficient is \eqref{eq:response}
 with both actual marked Hessians.
 \item The coefficient is invariant under compactly supported
 vertical gauge changes of any one of its three arguments
 supported in the tested interior plateau.
 It is strictly nonzero for real smooth finite packets with the
 prescribed normalization. Both Hessian terms are required in
 the Ward identity, even though the selected test is obtained
 by continuity from a limit in which they vanish.
 \item With the same centered physical Hadamard reference,
 \eqref{eq:response} is the normalized connected coefficient
 of one Wick-ordered tree vertex with a coherent incoming
 perturbation and the matched source/readout derivatives.
 Partitioning the same cubic vertex by a smooth accounting
 cut, retaining its source contacts, does not change the result.
\end{enumerate}
\end{theorem}

This construction uses the physical linearized Einstein--scalar
algebra of the kind studied in \cite{hack}, and a positive compact
Cauchy covariance from \cite{hadamard,modular}. The differentiated
Ward identity is the classical, one-vertex level of the usual
locally covariant framework \cite[Sec.~3.4]{bfr}. The additional
claim is the actual normalized mixed retarded test on one
active background, with its two genuine source contact terms.
No all-order Ward identity, ultraviolet coincidence theorem,
or independent subregion composition follows from it.

\section{Exact coupled stock and physical linear packets}

Let \((\Sigma,\gamma_*)\) be a closed orientable hyperbolic
three-manifold of sectional curvature \(-1\).
Use \(K_{ij}^{\rm ext}=-\tfrac12\Lie_n\gamma_{ij}\) for the
second fundamental form, distinct from the gauge differential
\(K\). At a marked slice with mean curvature \(\tau=-3\), take
\[
  \gamma=u^4\gamma_*,
  \qquad K^{\rm ext}=\tfrac{\tau}{3}\gamma,
  \qquad X=\epsilon f,\qquad \Pi_X=n(X)=0 ,
\]
where \(f\) is a fixed real smooth function nonconstant on a
small chosen patch. The exact Hamiltonian and momentum constraints
reduce to
\begin{equation}
 -8\Delta_*u-6u+6u^5-\epsilon^2|df|_{\gamma_*}^2u=0,
 \qquad 0=0 ,
 \label{eq:conformal}
\end{equation}
respectively. The derivative in \(u\) at \((u,\epsilon)=(1,0)\)
is \(-8\Delta_*+24\), with \(\Delta_*\le0\), an isomorphism.
The implicit function theorem gives a positive smooth
\(u=1+O(\epsilon^2)\). This is the elementary CMC conformal
construction for this stock \cite{cbip,icb}.
There are no conformal Killing fields on the compact negative
constant-curvature reference: the York operator is strictly
negative by the integrated identity using \(\operatorname{Ric}(\gamma_*)=-2\gamma_*\).
The property is conformally invariant. At the displayed cut the
tracefree first Killing-initial-data equation would make the
shift conformal Killing, hence zero. Its trace with \(\tau=-3\)
then makes the lapse zero. There is therefore no infinitesimal
Killing stabilizer on this stock.

The smooth Einstein--scalar Cauchy evolution gives a common
short slab after shrinking uniformly near \(\epsilon=0\);
the limiting solution is compact Milne,
\(\bar g_0=-\mathrm dt^2+t^2\gamma_*\) near \(t=1\), with
\(\bar X_0=0\). For the construction below, all intervals and
initial smooth profiles are chosen at finite order first, then
\(\epsilon>0\) is chosen sufficiently small. On the chosen
patch \(d\bar X\) is nonzero already on the initial surface
and remains nonzero in a short neighborhood.

We need actual tangents of the constraint surface. First vary
\(\gamma_*\) in a nonzero smooth TT tensor \(h_*\), keeping
\(\tau\) fixed, and solve the conformal equation again, now with
the varied scalar curvature \(\operatorname{Scal}(\gamma_*)\)
in place of \(-6\). This yields a metric-like Jacobi field \(U\).
At \(\epsilon=0\), the scalar-curvature derivative of a TT tensor
is zero, so the derivative of the conformal factor is zero.
The vacuum Milne tensor perturbation evolves in synchronous
spatial TT form; the spatial divergence and trace constraints
are preserved by the constant-curvature tensor wave operator.
In particular that perturbation is also in the harmonic wave
gauge used below. Nonzero smooth TT tensors are available, for
example as the linearized Cotton--York tensor of suitable compact
smooth metric variations on the constant-curvature reference
\cite{beig}; its TT identity follows by differentiating the
trace and divergence identities at the Cotton-flat metric.

Independently vary the scalar initial position by either
\(s_{1,0}\) or \(s_{2,0}\), leaving \(\Pi_X=0\); solve the same
scalar conformal constraint. This gives two smooth true
Jacobi fields \(s_1,s_2\). At \(\epsilon=0\) they are pure
scalar with zero initial normal derivative. For small nonzero
\(\epsilon\) they have the metric parts determined by the
coupled constraints and evolution. To verify that the selected
three fields are not null on the physical phase space, include
their symplectic mates: an independent TT shear-momentum
perturbation for \(U\), and independent scalar initial
velocity variations for \(s_1,s_2\). These can be completed
to constrained data. In the CMC conformal system the
linearized momentum return at the reference is the invertible
operator \(\Delta_*-2+\tfrac13 d\divg_*\) on one-forms.
At the reference maximal nonsingular splitting it has
\(-\|\nabla V\|^2-2\|V\|^2-\tfrac13\|\divg V\|^2\)
as its quadratic form. Together with the scalar block
\(-8\Delta_*+24\), it gives a split inverse for the constraints,
and hence the completed variations persist in the small
coupled branch. At \(\epsilon=0\), each of the three
canonical pairings is a nonzero multiple of the squared
norm of its real initial profile. The coupled reduced
pairings vary continuously, so they stay nonzero for small
\(\epsilon\).

\section{Exact closed currents and their Green representatives}

The linearized Noether relations are \(HK=0\) and \(K^*H=0\).
We use the harmonic wave gauge with
\[
 F(h,\sigma)_a
   =\nabla^b h_{ba}-\tfrac12\nabla_a\tr_{\bar g}h
                         -2\sigma\nabla_a\bar X .
\]
Let \(D=FK\) be its normally hyperbolic ghost operator.
For clarity about the field--density identification, use the
action pairings with trace-reversal on the metric field and
the scalar pairing, normalized so that
\(K^\dagger=-F/2\). For example the field pairing is
\(\tfrac14 h:I\tilde h+\sigma\tilde\sigma\),
\(Ih=h-\tfrac12\bar g\,\tr_{\bar g}h\), with the common
volume density. In this identification the wave-gauge Hessian
is \(P=H^\natural+KF\), where \(H^\natural\) denotes \(H\)
with its first slot raised using the action pairing.
Then \(PK=KD\) and \(FP=DF\). The Einstein--scalar
linearization cancels derivative mixing in \(P\) and gives
scalar wave principal symbols; these are the canonical
positive-kinetic wave-gauge identities of \cite{modular}.
They also follow directly by differentiating
\(\Ric=d\bar X\otimes d\bar X\), \(\Box\bar X=0\):
the tensor derivatives of \(\sigma\) cancel those in
\(K(-2\sigma\,d\bar X)\), and the scalar derivative of \(h\)
combines with \(F(h,0)(\bar X)\). Sources are lowered
or raised by this fixed real pairing, not by a choice of
normalization fitted to the output.

Take smooth time ramps \(b_-\) increasing from zero to one
and \(b_+\) decreasing from one to zero on two disjoint
compact subintervals, with the incoming ramp earlier.
Define the exact densities
\begin{equation}
 J_1=H(b_-s_1),\qquad J_2=H(b_+s_2).
 \label{eq:currents}
\end{equation}
Since \(Hs_i=0\), they are supported on their ramps;
compact \(\Sigma\) makes these spacetime-compact.
The exact identity \(K^*H=0\), not merely a principal
transversality test, proves they are closed. A scalar-only
function inserted in their place would usually fail this
condition when \(d\bar X\ne0\).

Writing \(G_D^R,G_D^A\) for the retarded and advanced inverses,
the actual one-sided wave-gauge representatives are
\begin{align}
 r&=b_-s_1-KG_D^R F(b_-s_1),\\
 a&=b_+s_2-KG_D^A F(b_+s_2).
 \label{eq:legs}
\end{align}
In these formulas the ghost inverses also act on functions
supported to one time side within the regular slab, by their
ordinary vanishing Cauchy conditions on the other side.
They make \(Fr=Fa=0\) and \(Hr=J_1,\ Ha=J_2\), with the
required one-sided support. In particular \(r\) and \(a\)
are not solutions of the homogeneous Hessian on their own
ramps. A compatible homogeneous gauge representative of
\(U\) is obtained by the same wave-gauge correction with
fixed initial gauge data. At the Milne endpoint the correction
is unnecessary for \(U\), and both scalar legs in \eqref{eq:legs}
are exactly \(b_-s_1,b_+s_2\).

For later normalization the same ramp realization applies to
any smooth physical homogeneous field \(v\) on the slab.
Choose \(Fv=0\) and a rising ramp \(b\). The retarded
and advanced solutions for the single current \(H(bv)\)
are built from \(bv\) and \((b-1)v\), respectively, plus
their one-sided gauge corrections. Their difference is
\(v\) modulo gauge. For the declining probe ramp the same
statement holds with the opposite sign. Thus the fields
and the two labelled packets have actual compact sources
on the physical Weyl phase space. Their positive
one-particle norms follow from the quasifree uncertainty
inequality and the nonzero symplectic pairings above.

\section{Geometric attachment and the two Ward contacts}

We now construct the source functional, so \(Q_J\) is fixed
before calculating the response. On a smaller common slab
label the CMC cuts by \(\tau(t)=-3/t\) near the marked slice.
For nearby smooth metric configurations, including off-shell
ones, first take their nearby CMC graph at each assigned
\(\tau(t)\). The graph derivative at Milne, with our choice
of future normal, is \(-\Delta_{\gamma(t)}+3/t^2\) (an
overall lapse sign has no effect on invertibility). Then fix
coordinates on that graph by the harmonic map to its
corresponding marked stock slice near the identity. The target
metric is the induced stock metric \(\bar\gamma_{\epsilon,t}\),
held fixed in all variations at this stock; the identity is
harmonic at the stock, also when \(\epsilon\ne0\). The targets
at \(t,\epsilon\) can be chosen as a smooth family using the
first CMC graph construction and a fixed near-identity spatial
identification. At Milne the map-variable derivative is the
invertible Jacobi operator given by a positive rescaling of
\(\Delta_*-2\), with \(\Delta_*\le0\). The normal equation
is solved first; the tangential dependence on that solution
does not enter the diagonal blocks of the combined derivative.

For the perturbed stock these inverses persist on
the compact interior family of cuts. The graph and harmonic
map equations are elliptic on each slice; the usual local
implicit-function construction gives smooth dependence for
the finite smooth metric families and Taylor directions used
here. Concretely one can first take a sufficiently high
Sobolev order on each cut and a finite number of derivatives
in the time parameter; the graph equation evaluates the
smooth nearby metric and its first derivatives, and its
first and second jets at a fixed cut depend only on the
corresponding jets of those metric data there. The harmonic
section likewise uses only the induced data on that cut.
We need this cutwise construction and two Taylor orders,
not smoothness of an undifferentiated same-order Sobolev
diffeomorphism group action. With the sign of the lapse
kept near the reference the ordered graphs form a smooth
slice family on a smaller slab. Elliptic uniqueness in the
chosen near-identity neighborhoods and identical markings
give covariance under sufficiently small diffeomorphisms
trivial near the fixed outer collars, since the geometric
conditions defining the slice and the marked map are
covariant. See \cite{am} for the CMC spatial-harmonic
elliptic gauges. Nonlinear relational pullbacks are
standard \cite{relational}; the particular
preparation/readout below use their two-jet.

Let \(\Psi(Y)\) denote the actual pullback of \(Y\) to
this marked chart, using markings fixed at the stock.
Thus \(\Psi(\bar Y)=\bar Y\) in the fixed reference
coordinates and \(\Psi(\varphi^*Y)=\Psi(Y)\) for the
local vertical gauge changes under discussion. Put
\[
 O_J(Y)=\langle J,\Psi(Y)-\bar Y\rangle .
\]
The density \(J\) is in the fixed marked chart. This is
not a bare tensor component with an external probe frozen
in a gauge-variant manner. The first displacement of the
pullback can be written \(-Z(v)\) for a variation \(v\),
with \(Z(K\xi)=\xi\) for interior gauge directions.
Thus \(O'_J=J\) at the stock by closure. Expanding the
actual pullback through second order gives
\begin{align}
 Q_J(v,w)=\big\langle J,\;&-\Lie_{Z(v)}w-\Lie_{Z(w)}v
 \nonumber\\
  &+\tfrac12
      (\Lie_{Z(v)}\Lie_{Z(w)}+\Lie_{Z(w)}\Lie_{Z(v)})
                                                \bar Y\big\rangle .
 \label{eq:hessian}
\end{align}
The second displacement generator multiplies \(K\) and
drops out against \(J\). This is just a Taylor expansion
of the pullback; it does not require the geometric
diffeomorphism to be the exponential of a time-independent
vector field. The chart may be spatially nonlocal on a
cut. For these Hessians the dependence on time remains in
the current's marked collar: both Taylor variations at a
given assigned cut use only the fields and a finite number
of their derivatives at that cut, with elliptic inverses
on the same spatial surface. In particular the Hessian
viewed as a source in its second slot is temporally compact
when \(J\) is. The measured supports are therefore
time collars on compact spatial surfaces, not claims about
spatially local apparatus.

For any allowed compactly supported gauge vector \(\xi\),
differentiate invariance of \(O_J\) and of the action with
all variations in their actual positions:
\begin{align}
 Q_J(K\xi,v)&=-\langle J,\Lie_\xi v\rangle ,
                                  \label{eq:wardsource}\\
 C(K\xi,v,w)&=-\langle Hv,\Lie_\xi w\rangle
                    -\langle Hw,\Lie_\xi v\rangle .
                                  \label{eq:wardaction}
\end{align}
In \eqref{eq:wardaction} the first variation of \(S\)
vanishes at the stock. No use of \(Hv=Hw=0\) is made
in forming this identity. In particular the two source
terms for \(v=r,w=a\) are both essential. The identities
remain exact in the tested compact time region; if one
inserts a fixed cubic cutoff, it is chosen to be one on
a neighborhood of the full compact overlap of \(r,a\),
including both ramps, and gauge changes are supported
inside the same plateau, with the same anchor.
The action derivatives with a tested compact gauge
argument then agree with those of the uncut
relative action. A cutoff changed where those variations
are supported would describe a different test.

\section{The sourced causal coefficient}

Prescribe the first homogeneous incoming variation
\(\alpha U\) and turn on \(-\lambda O_{J_1}\) in the
same action. The first retarded variation in \(\lambda\)
is \(r\) from \eqref{eq:legs}. At order
\(\alpha\lambda\), the field variation \(z\) solves the
retarded equation
\begin{equation}
 Hz=-C(U,r,\cdot)+Q_{J_1}(U,\cdot).
 \label{eq:z}
\end{equation}
The right side is itself a closed density:
\eqref{eq:wardaction} with \(HU=0,Hr=J_1\) and
\eqref{eq:wardsource} cancel on \(K\xi\).
This supplies a legitimate retarded Green prescription for
\(z\), modulo gauge. Only its part in the past of the probe
is used, so the one-sided retarded equation on the slab
suffices. Pairing with the probe by retarded--advanced
duality and differentiating the output \(O_{J_2}\) give
\begin{align}
 \partial_\alpha\partial_\lambda O_{J_2}\big|_0
 &=\langle J_2,z\rangle+Q_{J_2}(U,r)\nonumber\\
 &=-C(U,r,a)+Q_{J_1}(U,a)+Q_{J_2}(U,r).
 \label{eq:derivation}
\end{align}
All three terms are determined by the same action and
the same geometric attachment.

For \(U\mapsto U+K\xi\) the variation of the negative
cubic in \eqref{eq:derivation} is
\(\langle J_1,\Lie_\xi a\rangle+\langle J_2,\Lie_\xi r\rangle\).
The two Hessians cancel it separately.
For the change of \(r\) only the \(J_2\) term remains in
the cubic Ward relation (because \(HU=0\)), and the
change of \(Q_{J_2}\) cancels it. For \(a\), the identical
argument uses \(J_1\). These are changes of representatives
of the same fields in the tested support class; in a fixed
wave gauge the one-sided Green representatives themselves
are unique. In particular no assertion of a nontrivial
freedom to add a past-zero homogeneous wave-gauge ghost
to a retarded solution is needed.

These contact terms are not merely a formal option.
For example at the Milne endpoint on the early ramp,
\(J_1\) has scalar component \([\Box,b_-]s_1\)
while the advanced scalar leg is \(s_2\) there.
The time profiles can be chosen so that this component
is nonzero on an open set where \(\nabla s_2\) is
nonzero. A smooth spatial \(\xi\) supported inside
that set with suitable sign gives
\(\langle J_1,\Lie_\xi a\rangle\ne0\), and \(J_2\)
has no contribution there. The analogous late-ramp
choice isolates \(J_2\). For the real oscillatory initial
profile below, one can choose these open sets where
both the profile and its first spatial derivative are
nonzero, and choose the ramps to have a nonzero
second-derivative term there. These nonzero variations
persist for fixed choices at small positive \(\epsilon\).
The active scalar background may be chosen with
\(df\ne0\) in the common region. Thus the bare
uncompleted cubic does not define this representative-independent
test, even when the particular value of a Hessian
happens to vanish in a special gauge or at a limiting
background.

\section{An attained nonzero finite packet}

Select a point and a real coordinate direction \(x^1\)
where the contravariant component \(U^{11}\) of the
reference tensor mode is positive. Such a direction
exists because \(U\) is nonzero and trace-free.
Choose a small chart and nonzero real
\(\chi\in C_c^\infty\) supported where \(U^{11}>0\).
Take both initial scalar profiles
\(s_{1,0}=s_{2,0}=\chi\cos(Lx^1)\) and both
initial velocities zero. Then
\[
 \int_{\Sigma} U^{ij}\partial_i s_{1,0}
                 \partial_j s_{2,0}\,\dd\mu_{\gamma_*}
 =
 L^2\int_{\Sigma} U^{11}\chi^2\sin^2(Lx^1)
                 \,\dd\mu_{\gamma_*}+O(L).
\]
The leading integral tends to half the strictly positive
integral of \(U^{11}\chi^2\). Choose one large but finite
\(L\) so the displayed result is strictly positive,
then choose a short regular interval on which its evolved
version remains strictly positive. Take the ramps, with
\(0\le b_\pm\le1\), inside that interval, with nonzero
overlap \(b_-b_+\). The packet fields have their usual
finite physical supports on the Cauchy slab;
there is no plane wave or limit narrowing a detector
with \(\epsilon\).

At \(\epsilon=0\), the synchronous TT mode has zero
first displacement under the geometric marking:
its linear \(\tau\) is zero because it has zero
normal and trace perturbations and the reference
extrinsic curvature is pure trace; its spatial
harmonic tension vanishes by being TT on each cut.
The scalar legs at that endpoint have zero metric
variation. Hence \(Z(U)=Z(r)=Z(a)=0\) and both
Hessians in \eqref{eq:derivation} vanish there.
With the signs of \eqref{eq:action} the unique mixed
cubic is the derivative of the scalar kinetic term,
\[
 C(U,r,a)=\int_M
 \left(U^{ab}-\tfrac12\tr_{\bar g}(U)\bar g^{ab}\right)
             \partial_a r_X\,\partial_b a_X\,\dd\mu_{\bar g}.
\]
This \(hXX\) tensor vertex is familiar, for example
from inflationary cubic calculations \cite{maldacena}.
For the present finite test the trace and normal
components of \(U\) vanish and \(r_X=b_-s_1\),
\(a_X=b_+s_2\). The integral becomes
\[
 \int\! b_-(t)b_+(t)\dd t\int_{\Sigma_t}
 U^{ij}\partial_i s_1\,\partial_j s_2\,
                                   \dd\mu_{\gamma(t)}>0.
\]
There are no ramp derivatives left in this expression.
Consequently the full sourced coefficient is strictly
negative at the vacuum endpoint in our sign convention.

All amplitudes, causal solutions, gauge corrections,
and marked Hessians vary continuously in the sufficiently
smooth exact background family. Keeping the profiles and
all windows fixed gives \({\cal R}\ne0\) for
sufficiently small positive \(\epsilon\). One can choose
\(f\) and the chart with \(df\) nonzero on the initial
witness patch, so the resulting construction uses the
active scalar of the exact background, not a separate
spectator. The two scalar directions may be perturbed
independently if distinct profiles are wanted. The
nonzero result is at fixed finite settings of all
apertures and carrier frequencies in the perturbative
coefficient, not an asymptotic scattering claim.

\section{The normalized one-vertex coefficient}

Fix one of the allowed positive \(\epsilon\) backgrounds.
Use a centered positive quasifree Hadamard state on its
physical linear Einstein--scalar phase space, with
the physical Green CCR as in \cite{hadamard,modular}.
The scalar potential is zero and the Cauchy surface
compact, so the hypotheses of that construction are
satisfied. One may divide the three real packet
directions by their independently determined positive
one-particle covariance norms. For the two sourced
directions use the causal classes of the actual
\(J_1,J_2\); the realization following
\eqref{eq:legs} fixes their signs and supports.
Both criteria, finite norm and strict positivity,
hold on the selected smooth packet classes by
Hadamard covariance and their nonzero symplectic
mates. This changes \({\cal R}\) by a fixed nonzero
factor. No limit or continuity of these state norms
as \(\epsilon\) varies is needed.

For precision about which quantum coefficient is
meant, use one relative perturbation around the free
physical Gaussian algebra. Give its cubic vertex
a bookkeeper \(\kappa\), and Taylor-expand the same
marked source and probe at the same perturbative order:
\begin{align*}
 S_2(u)&+\frac{\kappa}{6}C(u,u,u)
      -\lambda\big(J_1(u)+\tfrac{\kappa}{2}Q_{J_1}(u,u)\big),\\
 {\cal O}_{2}(u)&=J_2(u)+\tfrac{\kappa}{2}Q_{J_2}(u,u).
\end{align*}
Constants and terms not contributing at order
\(\kappa\alpha\lambda\) are suppressed. A compact
physical source representing \(U\) implements a
unitary Weyl displacement; choose the Peierls sign
so its adjoint shifts the free field by \(\alpha U\),
and hence its displaced preparation is normalized.
The free linear drive of strength \(\lambda\)
shifts the retarded mean by \(\lambda r\).
Normal ordering is taken in the single fixed
Hadamard reference. At the fully connected
\(\kappa\alpha\lambda\) coefficient each surviving
slot is a smooth external packet or a causal leg:
one contraction with the linear readout takes the
last cubic leg to \(a\), and the two mean shifts
occupy the other two. The \(1/6\) is canceled by
the choice of readout leg and the permutation of the
two shifts. The quadratic marked source and probe
give the two cross terms, each with its \(1/2\)
canceled by symmetry. The coefficient is thus
exactly \eqref{eq:response}, with the stated
physical norms. This can equally be read as the
first connected retarded Bogoliubov coefficient of
the completed tested functional: its commutator
with a linear leg is the common causal propagator.
It uses no products of singular distributions
between two interaction vertices. We only assert
this finite tested tree multilinear form, not an
operator-valued extension of the entire relational
nonlinear functional in the free representation,
and certainly not positivity of an interacting
Hilbert-space representation. The closed density
in \eqref{eq:z} and the complete Ward cancellation,
rather than positivity on redundant metric
components, are what let the tested form descend
to the physical classes. Equivalently apply the linear marked projection
\(1-KZ\) to the three smooth external solutions on the
tested region. \(H\) is unchanged and \(Z\) annihilates
the projected representatives. The projection is local
in the cut label, so it preserves early/late vanishing;
its gauge generators can be smoothly cut off identically
one near the full tested overlap and both source supports.
Equation~\eqref{eq:hessian} then vanishes on every pair
of these marked directions, while Ward invariance gives
exactly the same coefficient as in wave gauge. This
expresses the smooth tested polynomial directly on the
positive physical solution range, without treating an
uncompleted redundant cubic as an operator.

Finally, write the same cubic profile as its
product with \(\theta\) plus with \(1-\theta\)
for a smooth accounting cut. The single vertex
is linear in its profile, so the two compactly
tested cubic integrals add to the original one.
The preparation and probe Hessians remain at
their original ramps (or can be partitioned as
the same distributions). If a split Ward
derivation differentiates the two complementary
switches, those terms cancel on addition.
This is an identity for one sourced functional
with fixed Green operator and fixed state;
it neither factors a transition into independently
prepared cut amplitudes nor changes a boundary
polarization.

\section{What the construction establishes}

The result is one nonzero mixed causal response of
ordinary Einstein--scalar dynamics on an attained
active background, computed with both source Hessians
and a common positive physical reference.
It tests a concrete interaction beyond the free
tensor and scalar covariance, but only at one vertex.
It does not supply quantum source backreaction,
separate gauge-constraint enforcement on the two
sides of a cut, closed interacting field algebras,
or a renormalized multivertex Ward identity.

\begin{thebibliography}{10}
\small\setlength{\itemsep}{1pt}

\bibitem{hack}
T.-P. Hack, Quantization of the linearised Einstein--Klein--Gordon
system on arbitrary backgrounds and the special case of perturbations
in inflation, \emph{Class. Quantum Grav.} \textbf{31}, 215004 (2014),
\href{https://arxiv.org/abs/1403.3957}{arXiv:1403.3957}.

\bibitem{hadamard}
Positive physical Hadamard states for critical Einstein--scalar
multiplets, research archive paper 09,
\href{https://github.com/g3173028019-beep/ai-automated-quantum-gravity-research/tree/main/papers/09-canonical-einstein-scalar-hadamard}{source and proof}.

\bibitem{modular}
Physical Hadamard states and bare modular kernels for coupled
Einstein--scalar systems, research archive companion 09+21,
\href{https://github.com/g3173028019-beep/ai-automated-quantum-gravity-research/tree/main/papers/09-21-coupled-modular-kernel}{source and proof}.

\bibitem{bfr}
R.~Brunetti, K.~Fredenhagen and K.~Rejzner,
Quantum gravity from the point of view of locally covariant quantum
field theory, \href{https://arxiv.org/abs/1306.1058}{arXiv:1306.1058}.

\bibitem{cbip}
Y.~Choquet-Bruhat, J.~Isenberg and D.~Pollack,
The constraint equations for the Einstein--scalar field system on
compact manifolds, \href{https://arxiv.org/abs/gr-qc/0610045}{arXiv:gr-qc/0610045}.

\bibitem{icb}
Y.~Choquet-Bruhat, J.~Isenberg and D.~Pollack,
Applications of theorems of Jean Leray to the Einstein--scalar
field equations,
\href{https://arxiv.org/abs/gr-qc/0611009}{arXiv:gr-qc/0611009}.

\bibitem{beig}
R.~Beig, TT-tensors and conformally flat structures on 3-manifolds,
\href{https://arxiv.org/abs/gr-qc/9606055}{arXiv:gr-qc/9606055}.

\bibitem{am}
L.~Andersson and V.~Moncrief,
Elliptic-hyperbolic systems and the Einstein equations,
\href{https://arxiv.org/abs/gr-qc/0110111}{arXiv:gr-qc/0110111}.

\bibitem{relational}
R.~Brunetti, K.~Fredenhagen, T.-P. Hack, N.~Pinamonti and K.~Rejzner,
Cosmological perturbation theory and quantum gravity,
\href{https://arxiv.org/abs/1605.02573}{arXiv:1605.02573}.

\bibitem{maldacena}
J.~Maldacena, Non-Gaussian features of primordial fluctuations in
single field inflationary models,
\href{https://arxiv.org/abs/astro-ph/0210603}{arXiv:astro-ph/0210603}.

\end{thebibliography}
\end{document}
